\section{Arquitectura y Protocolo de Sincronización}
\label{sec:arquitectura}

\subsection{Modelo de hilos}

El hilo principal (\texttt{main()}) actúa como el planificador; no existe un
hilo dedicado separado para esta responsabilidad. El programa simula un
único CPU lógico, por lo que no hay razón para que el planificador sea una
entidad concurrente distinta del programa que lo invoca: esta decisión
reduce en uno el número de hilos a crear, sincronizar y unir, y simplifica el
argumento de terminación. Además del hilo principal, se crea un
\texttt{pthread} por tarea (entre 5 y 25), cada uno ejecutando
\texttt{worker\_thread\_main}. Todo hilo creado se une con
\texttt{pthread\_join} antes de que el programa termine.

\subsection{Datos compartidos}

Dos estructuras concentran todo el estado que más de un hilo puede tocar:

\begin{itemize}
\item \textbf{\texttt{Task}} (una por tarea): identificador, boletos base
  (\texttt{tickets}), boletos efectivos (\texttt{effective\_tickets}, igual a
  \texttt{tickets} salvo compensación activa, Sección~\ref{sec:extension}),
  estado (\texttt{READY}/\texttt{RUNNING}/\texttt{FINISHED}), trabajo total y
  ejecutado, contador de despachos, estado privado de la serie de $\pi$, y su
  propia variable de condición \texttt{cond\_worker}.
\item \textbf{\texttt{Sync}} (una para todo el programa): un
  \texttt{pthread\_mutex\_t}, la variable de condición
  \texttt{cond\_scheduler} que espera el planificador, su predicado
  \texttt{event\_ready}, y la bandera \texttt{stop\_requested} que activa
  \texttt{--max-dispatches}.
\end{itemize}

Un único mutex (\texttt{sync.mutex}) protege el estado de todas las
\texttt{Task} (\texttt{state}, \texttt{dispatch\_count},
\texttt{effective\_tickets}, \texttt{debt}), \texttt{event\_ready} y
\texttt{stop\_requested}. Ese mutex \textbf{no} envuelve la ejecución de las
unidades de trabajo de $\pi$: hacerlo violaría la prohibición explícita del
enunciado de proteger toda la ejecución con un único bloqueo, y además sería
innecesario. La Sección~\ref{sec:invariante} explica por qué omitirlo ahí es
seguro.

\subsection{Variables de condición y predicados de espera}

Se usa un esquema de \textbf{una \texttt{cond\_worker} por tarea más una
\texttt{cond\_scheduler} compartida}, en vez de una única condvar global con
un predicado de ``turno'' señalada por \textit{broadcast}. La alternativa
descartada despertaría innecesariamente a las $N-1$ tareas que no ganaron
cada despacho --- el problema de \textit{thundering herd}: muchos hilos
despertados por un mismo evento cuando a lo sumo uno puede hacer trabajo
útil, desperdiciando ciclos de CPU y generando contención de mutex
innecesaria en los demás~\cite{MolloyLever2000}. El esquema elegido lo
evita señalando siempre de forma dirigida
(\texttt{pthread\_cond\_signal}, nunca \texttt{pthread\_cond\_broadcast}) a
un único hilo por evento. La Tabla~\ref{tab:condvars} resume las dos
condvars del protocolo: quién espera en cada una, su predicado exacto, y
quién la señala.

\begin{table*}[htbp]
\caption{Condvars, predicados y quién señala}
\label{tab:condvars}
\centering
\renewcommand{\arraystretch}{1.4}
\begin{tabular}{p{3.0cm}p{2.6cm}p{4.5cm}p{7.0cm}}
\toprule
\textbf{Condvar} & \textbf{Espera} & \textbf{Predicado} & \textbf{Señala} \\
\midrule
\texttt{task[i].cond\_worker} &
  Worker $i$ &
  \texttt{state==RUNNING || stop\_requested} &
  El planificador, al elegir a $i$ como ganador (o \texttt{sync\_request\_stop} al cortar) \\
\midrule
\texttt{cond\_scheduler} &
  El planificador &
  \texttt{event\_ready} &
  El worker en ejecución, al ceder o terminar \\
\bottomrule
\end{tabular}
\end{table*}

Ambos predicados se evalúan dentro de un \texttt{while}, nunca un
\texttt{if}, por dos razones independientes: POSIX permite explícitamente
que \texttt{pthread\_cond\_wait} retorne sin que nadie haya señalado la
condición~\cite{IEEE1003.1-2024} --- un despertar sin causa aparente,
conocido en la literatura de \textit{pthreads} como \textit{spurious
wakeup}~\cite{Butenhof1997}; y el caso en que varios hilos comparten una
condvar y el que despierta encuentra que la condición ya no aplica. \texttt{pthread\_cond\_wait} suelta el mutex y duerme el hilo de
forma atómica, y vuelve a tomar el mutex antes de retornar --- evita la
condición de \textit{señal perdida} que ocurriría si soltar el mutex y
dormirse fueran dos pasos separados.

\subsection{Protocolo de despacho}
\label{sec:protocolo}

La Fig.~\ref{fig:protocolo} muestra el ciclo completo de un despacho (el
pseudocódigo correspondiente es el Algoritmo~\ref{alg:bucle},
Sección~\ref{sec:algoritmos}). El planificador (hilo principal) ejecuta
tres pasos en secuencia:
\texttt{sync\_select\_winner} (sortea la ganadora entre las tareas
\texttt{READY}, ponderando por \texttt{effective\_tickets}),
\texttt{sync\_dispatch} (marca a la ganadora \texttt{RUNNING} y la señala), y
\texttt{sync\_wait\_for\_event} (se bloquea hasta que esa tarea ceda o
termine). El worker ganador, mientras tanto, estaba bloqueado en
\texttt{sync\_wait\_for\_dispatch}; al despertar, ejecuta su porción de
trabajo \textbf{fuera del mutex} y cierra el ciclo con
\texttt{sync\_finish\_turn}, que actualiza su estado y despierta al
planificador.

\begin{figure}[htbp]
\centering
\begin{tikzpicture}[
    box/.style={draw=black, rounded corners=3pt, text width=3.65cm,
                minimum height=0.9cm, align=center, font=\scriptsize},
    arr/.style={-{Stealth[length=2mm]}, thick},
    lbl/.style={font=\tiny, align=center},
    node distance=6mm
]
\node[box, fill=blue!8] (sel) {\texttt{sync\_select\_winner}\\(toma mutex, suma \texttt{effective\_tickets} de READY, sortea, suelta mutex)};
\node[box, fill=blue!8, below=of sel] (disp) {\texttt{sync\_dispatch}\\(toma mutex, assert exclusión, \texttt{state=RUNNING}, \texttt{signal(cond\_worker[i])}, suelta mutex)};
\node[box, fill=blue!8, below=of disp] (wait) {\texttt{sync\_wait\_for\_event}\\(toma mutex, \texttt{while(!event\_ready) cond\_wait}, suelta mutex)};

\node[box, fill=orange!12, below=16mm of wait] (wd) {\texttt{sync\_wait\_for\_dispatch}\\(worker $i$, dormido hasta \texttt{state==RUNNING})};
\node[box, fill=green!12, below=of wd] (work) {\texttt{workload\_run\_units}\\ \textbf{(sin mutex)} --- seguro por el invariante de exclusión};
\node[box, fill=orange!12, below=of work] (fin) {\texttt{sync\_finish\_turn}\\(toma mutex, \texttt{state=READY|FINISHED}, \texttt{event\_ready=1}, \texttt{signal(cond\_scheduler)}, suelta mutex)};

\draw[arr] (sel.south) -- (disp.north);
\draw[arr] (disp.south) -- (wait.north);
\draw[arr] (disp.east) to[out=-20,in=90] node[lbl, right=1pt] {signal\\dirigida} (wd.north east);
\draw[arr] (wd.south) -- (work.north);
\draw[arr] (work.south) -- (fin.north);
\draw[arr] (fin.west) to[out=160,in=-90] node[lbl, left=1pt] {signal} (wait.south west);
\draw[arr, dashed] (fin.east) to[out=20,in=-90] node[lbl, right=1pt] {si READY,\\vuelve a\\esperar} (wd.south east);
\end{tikzpicture}
\caption{Un ciclo de despacho. Las cajas azules corren en el hilo
planificador (\texttt{main}); las naranjas y la verde, en el hilo worker de
la tarea ganadora. Solo el tramo verde corre sin el mutex tomado.}
\label{fig:protocolo}
\end{figure}

\subsection{Por qué es seguro ejecutar el trabajo sin el mutex}
\label{sec:invariante}

El invariante de exclusión --- como máximo una tarea en \texttt{RUNNING} en
todo momento, garantizado por el \texttt{assert} dentro de
\texttt{sync\_dispatch} --- implica que mientras una tarea corre, ningún
otro hilo lee ni escribe sus campos: el planificador está bloqueado en
\texttt{sync\_wait\_for\_event} (no toca ninguna \texttt{Task}), y los demás
workers están bloqueados en \texttt{sync\_wait\_for\_dispatch} esperando su
propia señal (nunca tocan una tarea ajena). El único hilo que puede tocar la
tarea mientras corre es su propio dueño. Esta es una invariante del
\emph{protocolo} (el orden de llamadas), no algo que un \textit{sanitizer}
pueda verificar por sí solo --- de ahí que cualquier código nuevo que
necesite leer el progreso de una tarea \texttt{RUNNING} desde otro hilo (por
ejemplo, un log en vivo) deba traer esa lectura dentro del mutex
explícitamente, en vez de asumir que es seguro por analogía con este caso.

\subsection{Terminación limpia y \texttt{--max-dispatches}}

El enunciado exige que ningún trabajador quede bloqueado al finalizar. Una
tarea que llega a \texttt{FINISHED} simplemente retorna de
\texttt{worker\_thread\_main}. El caso menos obvio es
\texttt{--max-dispatches}: cuando el planificador alcanza el límite, una
tarea todavía \texttt{READY} que nunca vuelve a ser elegida se quedaría
esperando su señal para siempre. \texttt{sync\_request\_stop} resuelve esto
señalando la \texttt{cond\_worker} de cada tarea \texttt{READY} sin cambiar
su estado; \texttt{sync\_wait\_for\_dispatch} retorna un código distinto de
cero en ese caso, y el worker retorna de inmediato sin ejecutar trabajo ni
llamar \texttt{sync\_finish\_turn}, dejando su tarea intacta en
\texttt{READY}. El planificador entonces une los $N$ hilos con
\texttt{worker\_pool\_join} sin que ninguno quede colgado.
